<?xml version="1.0" encoding="UTF-8"?>
<worksheet>
  <title>Generalized eigenspaces</title>

  <introduction>
    <p>
      In this worksheet, we explore generalized eigenspaces and how they
      decompose a vector space. Let <m>T:V \to V</m> be a linear operator
      with minimal polynomial <m>p(x)=(x-\lambda)^k q(x)</m>, where
      <m>q(\lambda)\neq 0</m>. Then <m>\lambda</m> is an eigenvalue and
    </p>
    <p>
      <me>G_\lambda = \nul((T-\lambda I)^k)</me>
    </p>
    <p>
      is called the generalized eigenspace. Elements of <m>G_\lambda</m>
      are called generalized eigenvectors.
    </p>
  </introduction>

  <exercise>
    <task workspace="1in">
      Suppose that <m>\vvec</m> is a generalized eigenvector with
      eigenvalue <m>\lambda</m>. Explain why
      <m>(T-\lambda I)^k\vvec = 0</m>.
    </task>

    <task workspace="1in">
      What can you conclude about the annihilating polynomial
      <m>p_\vvec</m>?
    </task>

    <task workspace="1in">
      Suppose that <m>\lambda</m> and <m>\mu</m> are two distinct eigenvalues.
      Use these observations to explain why
      <m>G_\lambda \cap G_\mu = \{0\}</m>.
    </task>
  </exercise>

  <exercise>
    <task workspace="1in">
      Suppose that <m>\dim V = 4</m> and <m>V_1</m> and <m>V_2</m> are subspaces
      of <m>V</m> with <m>\dim V_1 = 3</m> and <m>\dim V_2 = 2</m>. Then there
      are bases <m>\{\vvec_1,\vvec_2,\vvec_3\}</m> for <m>V_1</m> and
      <m>\{\vvec_4,\vvec_5\}</m> for <m>V_2</m>. Explain why the set
    </task>

    <task workspace="1in">
      <p>
        <me>\{\vvec_1,\vvec_2,\vvec_3,\vvec_4,\vvec_5\}</me>
      </p>
    </task>

    <task workspace="1in">
      is linearly dependent.
    </task>

    <task workspace="1in">
      Use this fact to show that
      <m>V_1 \cap V_2 \neq \{0\}</m>.
    </task>
  </exercise>

  <exercise>
    <task workspace="1in">
      Suppose that <m>p(x)=(x-2)^3(x-4)^2</m>. If <m>\wvec</m> is in
      <m>\range((T-4I)^2)</m>, explain why
      <m>(T-2I)^3(\wvec) = 0</m>.
    </task>

    <task workspace="1in">
      Explain why
      <m>\range((T-4I)^2) \subset \nul((T-2I)^3) = G_2</m> and hence
      <m>\dim \range((T-4I)^2) \leq \dim G_2</m>.
    </task>

    <task workspace="1in">
      Explain why
      <p>
        <me>\dim G_4 + \dim \range((T-4I)^2) = \dim V.</me>
      </p>
    </task>

    <task workspace="1in">
      Explain why
      <p>
        <me>\dim G_4 + \dim G_2 \geq \dim V.</me>
      </p>
    </task>

    <task workspace="1in">
      Since <m>G_4 \cap G_2 = \{0\}</m>, explain why
      <m>\dim G_4 + \dim G_2 = \dim V</m>.
    </task>

    <task workspace="1in">
      Explain why <m>G_4 \oplus G_2 = V</m>.
    </task>

    <task workspace="1in">
      Explain why both <m>G_4</m> and <m>G_2</m> are <m>T</m>-invariant.
    </task>

    <task workspace="1in">
      If <m>\bcal</m> is a basis for <m>G_4</m> and <m>\ccal</m> is a basis for
      <m>G_2</m>, then <m>\dcal=\bcal\cup\ccal</m> is a basis for <m>V</m>.
      Explain why
    </task>

    <task workspace="1in">
      <p>
        <me>
        \coords{T}{\dcal,\dcal} =
        \begin{bmatrix}
        \coords{T}{\bcal,\bcal} &amp; \zerovec \\
        \zerovec &amp; \coords{T}{\ccal,\ccal}
        \end{bmatrix}.
        </me>
      </p>
    </task>
  </exercise>

  <exercise>
    <task workspace="1in">
      Suppose the minimal polynomial factors as
    </task>

    <task workspace="1in">
      <p>
        <me>
        p(x) =
        (x-\lambda_1)^{k_1}(x-\lambda_2)^{k_2}\ldots(x-\lambda_m)^{k_m}.
        </me>
      </p>
    </task>

    <task workspace="1in">
      Explain why
    </task>

    <task workspace="1in">
      <p>
        <me>
        V = G_{\lambda_1} \oplus G_{\lambda_2} \oplus \ldots \oplus
        G_{\lambda_m}.
        </me>
      </p>
    </task>

    <task workspace="1in">
      Explain why there is a basis for <m>V</m> such that the matrix
      representing <m>T</m> is block diagonal with blocks corresponding
      to <m>T|_{G_{\lambda_j}}</m>.
    </task>
  </exercise>

  <conclusion>
    <p>
      In this worksheet, you have shown that generalized eigenspaces corresponding
      to distinct eigenvalues intersect trivially and together span the vector
      space. This leads to a decomposition of <m>V</m> into invariant subspaces,
      providing a structural understanding of linear operators via their minimal
      polynomials.
    </p>
  </conclusion>

